\subsection{Automated Synthetic Environment Generation Pipeline}

We propose an automated pipeline that converts trajectory-level failures in the
target environment into capability-specific training environments. The pipeline
consists of two stages: \emph{contrastive skill identification} and
\emph{environment synthesis}.

\paragraph{Contrastive Skill Identification.}
Given a dataset $\mathcal{D} = \{(\tau_i, r_i)\}_{i=1}^{N}$ of agent
trajectories $\tau_i$ and their binary outcomes $r_i \in \{0, 1\}$ from the
target environment, we analyze \emph{both} successful trajectories
$\mathcal{D}^{+} = \{(\tau_i, r_i) \mid r_i = 1\}$ and failed trajectories
$\mathcal{D}^{-} = \{(\tau_i, r_i) \mid r_i = 0\}$ to identify capability
gaps. For each trajectory, we extract the sequence of tool calls, tool results,
and the agent's final response, then annotate which capabilities were exercised
and whether each was executed successfully. We operationalize this annotation
using a combination of programmatic checks against ground-truth action
sequences and LLM-assisted classification for ambiguous cases.

This produces a set of candidate capabilities $\mathcal{C}$. For each
$c \in \mathcal{C}$, we estimate two complementary statistics:
\begin{align}
    \text{SR}^{+}(c) &= P(\text{$c$ executed correctly} \mid
        \text{$c$ required}, \, \tau \in \mathcal{D}^{+}), \\
    \text{SR}^{-}(c) &= P(\text{$c$ executed correctly} \mid
        \text{$c$ required}, \, \tau \in \mathcal{D}^{-}).
\end{align}
$\text{SR}^{+}(c)$ measures how reliably capability $c$ is executed in
successful trajectories, while $\text{SR}^{-}(c)$ measures the same in failed
trajectories. We select capabilities with a large gap between these rates:
\begin{equation}
    \mathcal{C}^{*} = \bigl\{c \in \mathcal{C} \;\big|\;
        \text{SR}^{+}(c) - \text{SR}^{-}(c) \geq \delta \bigr\},
\end{equation}
retaining capabilities that are consistently executed in successful trajectories
but frequently fail in unsuccessful ones. This contrastive criterion is more
robust than analyzing failures alone: it filters out capabilities that have
universally low success rates (indicating inherent task ambiguity or annotation
noise rather than a trainable deficit) and capabilities that fail in isolation
but do not actually discriminate between successful and unsuccessful
trajectories. To further improve robustness, we perform the annotation process
across multiple independent passes and retain only capabilities that are
consistently identified.

\paragraph{Environment Synthesis.}
For each retained capability $c \in \mathcal{C}^{*}$, we construct a synthetic
micro-environment $\mathcal{E}_c = (\mathcal{T}_c, \mathcal{A}, R_c)$
consisting of a procedural task distribution $\mathcal{T}_c$, a shared action
space $\mathcal{A}$, and a programmatic reward function $R_c$.

\textbf{Shared action space.}
Each micro-environment preserves the target environment's tool schemas,
multi-turn conversational protocol, and response format. That is, the agent
interacts with the same set of API tools (e.g., \texttt{get\_user\_details},
\texttt{cancel\_reservation}) using identical function signatures and system
instructions. This shared action space is the primary mechanism through which
training gains transfer to the target environment: capabilities acquired during
training compose directly with the target task distribution without additional
adaptation.

\textbf{Procedural task generation.}
The task distribution $\mathcal{T}_c$ is implemented as a seeded procedural
generator that produces diverse scenario instances. Given a random seed $s$,
the generator deterministically constructs a complete episode specification:
synthetic user profiles, database records, and task-specific parameters such as
constraint configurations and difficulty levels. Procedural generation serves
two purposes: (i)~it prevents the agent from memorizing solutions, since each
seed yields a unique scenario, and (ii)~it controls difficulty and diversity
through configurable parameters (e.g., number of database records, number of
constraints, number of sequential operations).

\textbf{Capability-isolating reward.}
The reward function $R_c$ is designed so that success depends primarily on
correct execution of capability $c$, while minimizing dependence on unrelated
capabilities. We implement $R_c$ as a programmatic verifier that checks the
agent's actions against the ground-truth solution---for example, comparing the
resulting database state hash against the expected state, or verifying that
specific tool calls were made with correct arguments. This yields a dense,
automatically computed reward signal without requiring human annotation or
LLM-based evaluation.

The resulting set of micro-environments
$\{\mathcal{E}_c\}_{c \in \mathcal{C}^{*}}$ provides targeted, dense
supervision for each identified capability gap. Because the environments are
modular, we compose them into a training curriculum by sampling environments
with tunable mixing weights $\{w_c\}_{c \in \mathcal{C}^{*}}$, concentrating
training compute on the capabilities with the largest performance gaps.
